\scanpage{18}
\mspage{- 2 -}
\figure{\begin{tikzpicture}[x=1mm,y=1mm]
\fill (0,2.5)--(1.5,3.1)--(4.1,2.8)--(5.5,2.1)--(7.3,1.9)--(9.2,.6)--(8.7,-.3)--(6.1,.15)--(2.5,.5)--cycle;
\end{tikzpicture}}

the surjectivity $B\twoheadrightarrow A$ will induce the antipodal map on the $k$-algebra $A$ and thus ``force'' it to be a group object (not the more general monoid object) in $\mathrm{Aff}_k$.

Consider the example
\[
\begin{aligned}
\{0,1\}: \mathrm{Alg}_k&\longrightarrow\mathrm{Sets}\\
k'&\longmapsto\mathrm{Hom}_{k\text{-alg}.}(k\times k,k').
\end{aligned}
\]

Make $A=k\times k$ into a coalgebra over $k$ by
\[
\begin{aligned}
 A&\xrightarrow{\unsure{\Pi_0|\Pi_o}}A\otimes A\\
 (a,b)&\longmapsto.\ (a,b)\otimes(1,0)+(1,1)\otimes(0,b).
\end{aligned}
\]

the two projections.
\[
\begin{aligned}
P_1 &: (a,b)\longmapsto a.\\
P_2 &: (a,b)\longmapsto b
\end{aligned}
\]
represent $k$-valued points of $\{0,1\}$

multiplication table :
\diagram{$\begin{array}{c|cc}
\Pi & P_1 & P_2\\ \hline
P_1 & \multicolumn{1}{c|}{P_1} & P_2\\ \cline{2-3}
P_2 & \multicolumn{1}{c|}{P_2} & P_2
\end{array}$}

(this is easy to check :
\[
\begin{aligned}
(P_1,P_1)&:(a,b)\otimes(c,d)&\longmapsto\unsure{ac|\acute a\acute c|a'c'}\unsure{.}\\
(P_1,P_2)&:\qquad\text{\rq\rq}&\longmapsto ad\\
(P_2,P_1)&:\qquad\text{\rq\rq}&\longmapsto bc\\
(P_2,P_2)&:\qquad\text{\rq\rq}&\longmapsto bd.
\end{aligned}
\]
% archive sheet number: 17
\clearpage
